% Generated draft: requires mathematical and source review.
\documentclass[11pt]{article}
\usepackage[margin=25mm]{geometry}
\usepackage{fontspec}
\setmainfont{Latin Modern Roman}
\usepackage{amsmath,amssymb,mathtools,unicode-math}
\setmathfont{Latin Modern Math}
\usepackage{enumitem,needspace,xurl,hyperref}
\setlength{\parindent}{0pt}
\setlength{\parskip}{5pt}
\setlength{\emergencystretch}{3em}
\newcommand{\N}{\mathbb{N}}
\newcommand{\Z}{\mathbb{Z}}
\newcommand{\Q}{\mathbb{Q}}
\newcommand{\R}{\mathbb{R}}
\newcommand{\C}{\mathbb{C}}
\newcommand{\F}{\mathbb{F}}
\newcommand{\A}{\mathbb{A}}
\newcommand{\PP}{\mathbb{P}}
\newcommand{\E}{\mathbb{E}}
\newcommand{\eps}{\varepsilon}
\newcommand{\dd}{\,\mathrm{d}}
\begin{document}
\section*{UnsolvedMath ID: 30002453}
\textbf{Source problem code:} catalogue entry
\subsection{Definitions and mathematical statement}
Let \(k\) be a field and let \(\Lambda\) be a basic finite-dimensional \(k\)-algebra. The supplied statement imposes no further condition on \(k\). For notation, use finitely generated right \(\Lambda\)-modules, and let \(\tau\) denote the Auslander--Reiten translate. Write \(|X|\) for the number of isomorphism classes of indecomposable direct summands of a module \(X\), and let \(n\) be the number of isomorphism classes of simple \(\Lambda\)-modules.

A basic support \(\tau\)-tilting pair is a pair \((M,P)\), with \(M\) a basic module and \(P\) a basic projective module, satisfying
\[
\operatorname{Hom}_{\Lambda}(M,\tau M)=0,\qquad
\operatorname{Hom}_{\Lambda}(P,M)=0,\qquad
|M|+|P|=n.
\]
The set \(\mathrm{s}\tau\text{-}\mathrm{tilt}\,\Lambda\) consists of the isomorphism classes of basic support \(\tau\)-tilting modules, equivalently their associated basic support \(\tau\)-tilting pairs. For a module \(M\), let \(\operatorname{Fac}M\) be the full subcategory of factor modules of finite direct sums of copies of \(M\). Use the conventional order
\[
M\geq N\quad\Longleftrightarrow\quad
\operatorname{Fac}M\supseteq\operatorname{Fac}N.
\]

Let \(K_0(\operatorname{proj}\Lambda)\) be the split Grothendieck group of finitely generated projective modules. Choosing an ordering of the indecomposable projective modules identifies this group with \(\Z^n\). Given a minimal projective presentation
\[
P_1^M\longrightarrow P_0^M\longrightarrow M\longrightarrow0,
\]
define
\[
g(M)=[P_0^M]-[P_1^M],\qquad
g(M,P)=g(M)-[P].
\]
To distinguish these vectors from the cone data mentioned in the supplied background, write \(M=\bigoplus_{i=1}^{r}M_i\) and \(P=\bigoplus_{j=1}^{s}P_j\) as sums of indecomposable modules, and define the associated cone by
\[
C(M,P)=\left\{
\sum_{i=1}^{r}a_i g(M_i)-\sum_{j=1}^{s}b_j[P_j]
:\ a_i,b_j\in\R_{\geq0}
\right\}
\subseteq K_0(\operatorname{proj}\Lambda)\otimes_{\Z}\R.
\]

The supplied question is: for every basic finite-dimensional \(k\)-algebra \(\Lambda\), do the \(g\)-vectors determine the partial order on \(\mathrm{s}\tau\text{-}\mathrm{tilt}\,\Lambda\)? In the order convention above, the intended recovery must determine, for each pair of support \(\tau\)-tilting objects, whether \(\operatorname{Fac}M\supseteq\operatorname{Fac}N\).

There is an unresolved qualification: the question does not specify whether the retained data are module vectors, pair vectors, vectors of individual summands grouped by support \(\tau\)-tilting pair, or the associated cones. Nor does it specify precisely what labels, projective basis, or ambient lattice structure are retained in the word ``determine.'' Consequently, the source does not specify a unique reconstruction assertion. The cone definition above explains the background terminology without identifying these different possible inputs.

An algebra is called \(\tau\)-tilting finite here if \(\mathrm{s}\tau\text{-}\mathrm{tilt}\,\Lambda\) is finite. The supplied dated assessment labels the problem partially solved and states that the \(g\)-vector cones determine the order in this finite case, while the arbitrary finite-dimensional case remains conjectural. This is supplied curation, not an independently verified theorem or current literature assessment.
\subsection{Short English statement}
For a basic finite-dimensional algebra, the problem asks whether its support \(\tau\)-tilting order can be recovered from \(g\)-vector information. The supplied background reports a result for \(\tau\)-tilting finite algebras, but the exact data intended by ``\(g\)-vectors determine'' remain unspecified. This is a statement-only draft requiring human mathematical review.
\subsection{Sources}
\begin{itemize}
\item Local UnsolvedMath registry, permanent ID \(30002453\), source problem code \texttt{catalogue entry}, title: Recovery of Support \(\tau\)-Tilting Orders from \(g\)-Vectors. The supplied record belongs to the Oberwolfach Reports Open Problems collection. No individual proposer is identified.
\item Supplied report citation: \emph{Cluster Algebras and Related Topics} (2014), Oberwolfach Reports, pp.~3379--3432, DOI \texttt{10.4171/owr/2013/58}. \url{https://doi.org/10.4171/owr/2013/58}. The background describes Oberwolfach Reports as documenting mathematical workshops and their open-problem discussions. The record separately gives the proposed year as 2013.
\item Supplied literature source, labeled ``tau-tilting finite algebras and g-vectors'': \url{https://arxiv.org/abs/1503.00285}. No authors or fuller bibliographic details are supplied for this entry.
\item Literature assessment dated exactly \texttt{2026-08-21}: the supplied curation labels the problem partially solved and reports determination by \(g\)-vector cones for \(\tau\)-tilting finite algebras, with the arbitrary finite-dimensional case remaining conjectural. This assessment is supplied dated curation, not independently verified. The supplied background explicitly requests independent verification of time-sensitive status and cited claims.
\item The supplied registry labels the statement \texttt{corrected\_verified} and reports that it was checked against the official Oberwolfach report. These are supplied provenance claims; no independent source visit or verification is asserted here.
\end{itemize}
\end{document}
